\chapter{Fazit}
\label{chap:fazit}

Zwei Fragen hat diese Arbeit verfolgt.
Die erste galt der Gütegarantie von MULTIFIT und der Herleitung, welche zu ihr führt.
Die zweite galt der Subroutine, also dem Bin-Packing-Verfahren, das über jede getestete
Kapazität entscheidet.
Beide Stränge treffen sich in Implikation~\eqref{eq:ffd-akzeptiert}: Dort wird aus einer
Eigenschaft der Subroutine eine Aussage über den Makespan.

\paragraph{Die Performanzanalyse.}
Kapitel~\ref{chap:performanz} führt den Nachweis $r_m \le 5/4$ vollständig aus
(Theorem~\ref{thm:schranke54}).
Getragen wird er von Struktureigenschaften eines minimalen Gegenbeispiels: dem Löschlemma, dem
Dominanzlemma, der Mindestzahl von Aufgaben je Bin der $q$-Packung und den Größenschranken.
Theorem~\ref{thm:guete} verbindet das Grenzverhältnis anschließend mit dem Verfahren und liefert
$R_m(\mathrm{MF}(k)) \le r_m + 2^{-k}$.
Die Binärsuche kostet also nur einen additiven Term, welcher sich mit jedem Schritt halbiert.
In die Gegenrichtung führt Lemma~\ref{lem:untere-schranke} den Nachweis $r_m \ge 13/11$ für
$m \ge 13$ an einer Instanz von \textcite{friesen1984}.
Zusammen mit der oberen Schranke von \textcite{cao1995} ist das Grenzverhältnis für diese
Maschinenzahlen exakt bestimmt.

Die Schranke $13/11$ selbst bleibt in dieser Arbeit unbewiesen; Abschnitt~\ref{sec:beweisidee1311} stellt nur
ihre Beweisidee dar.
Caos Fallunterscheidung über $\Delta$ füllt sechzehn Seiten und stützt sich zusätzlich auf
\textcite{yue1990}.
Zwischen der hier bewiesenen Schranke $5/4$ und dem für $m \ge 13$ exakten Wert $13/11$ bleibt
somit eine Lücke; geschlossen wird sie in \cite{cao1995} über diese Fallunterscheidung.

\paragraph{Der Austausch der Subroutine.}
Kapitel~\ref{kap:alternative} prüft drei Alternativen, und die Befunde fallen verschieden aus.
Für MFFD gibt Lemma~\ref{lem:kollaps} eine Klasse von Instanzen an, auf der sich die Wahl
überhaupt nicht auswirkt: Ist die mittlere Maschinenlast mindestens doppelt so groß wie die
größte Aufgabe, in Kurzform $\rho(\Gamma) \ge 2$, so erzeugen MFFD und FFD bei jeder
getesteten Kapazität dieselbe Packung.
Geprüft ist das Kriterium in einem einzigen Durchlauf der Aufgabenliste.
Beim $\mathrm{OPT}+1$-Algorithmus von \textcite{jansen_solisoba2010} entsteht durch die
Modifikation zu einem $\varepsilon$-dualen Verfahren eine eigene Garantie:
Korollar~\ref{kor:multifit-jansen-dual} liefert den Makespan
$(1+\varepsilon)(1+2^{-k}) \cdot C^{*}_{\max}(\Gamma)$ mit $\varepsilon = 1/(2d-1)$, wobei $d$
die Anzahl verschiedener Aufgabengrößen bezeichnet.
Neben der Schranke für MULTIFIT+FFD ist sie die einzige instanzunabhängige Makespan-Garantie
dieser Arbeit; für den $\mathrm{OPT}+1$-Algorithmus selbst gibt Abschnitt~\ref{sub:anwendbarkeit}
eine von der Instanz abhängige Schranke an.

Für das Konfigurations-LP von \textcite{gilmore_gomory1961} fällt das Ergebnis negativ aus.
Lemma~\ref{lem:gg-rundung} lässt nach dem Runden $m + d - 1$ Bins zu und damit gerade die
$d-1$ Bins, die das Akzeptanzkriterium nicht mehr zulässt;
Implikation~\eqref{eq:ffd-akzeptiert} ist folglich nicht erfüllt.
Was das Verfahren stattdessen mitbringt, ist eine zertifizierte Ablehnung: Aus
$\lceil \mathrm{LP}[\Gamma,C] \rceil > m$ folgt $C < C^{*}_{\max}(\Gamma)$, während FFD im
entscheidenden Fall $\mathrm{FFD}[\Gamma,C] = m+1$ nichts über
$\mathrm{OPT}[\Gamma,C]$ aussagt.

\paragraph{Was die Messungen zeigen.}
Auf den \EvalInstanzen{} Instanzen der fünf Benchmark-Datensätze überschreitet kein gemessenes
Verhältnis $13/11$; der größte Wert beträgt \EvalMaxRatioOpt.
MULTIFIT+MFFD liefert dabei in \EvalWins{} Fällen einen kleineren und in \EvalLosses{} Fällen
einen größeren Makespan als MULTIFIT+FFD.
Sämtliche \EvalAbweichungen{} Abweichungen liegen unterhalb von $\rho = 2$, wie es
Lemma~\ref{lem:kollaps} verlangt.
Ein systematischer Vorteil der aufwendigeren Subroutine zeigt sich auf diesen Instanzen nicht.

Die zweite Versuchsreihe trennt Bin-Packing-Güte und Makespan noch deutlicher.
Auf dem HM-Gitter erreicht MULTIFIT+FFD auf \HmOptimalffd{} der Instanzen den optimalen
Makespan, mit dem $\mathrm{OPT}+1$-Algorithmus sind es \HmOptimaljansen.
Gegenüber der Grundlinie gewinnt MULTIFIT mit diesem Algorithmus dennoch nicht durchgängig:
\HmJansenBesserFFD{} Instanzen fallen besser aus, \HmJansenSchlechterFFD{} schlechter.
Den Grund liefert Abschnitt~\ref{sub:hm-garantien}: FFD verwirft \HmZuUnrechtFfd{} Kapazitäten,
für die eine Packung in $m$ Bins existiert, der $\mathrm{OPT}+1$-Algorithmus \HmZuUnrechtJansen{}
- nur an anderen Stellen.
Für den Makespan zählt demnach nicht, wie nah eine Subroutine an $\mathrm{OPT}[\Gamma,C]$
herankommt, sondern welche Kapazitäten sie akzeptiert; genau an diese Entscheidung knüpft
Implikation~\eqref{eq:ffd-akzeptiert} die Gütegarantie von MULTIFIT+FFD.
Die Entscheidung allein trägt den Makespan jedoch nicht: Der $\varepsilon$-duale Algorithmus
akzeptiert nach Abschnitt~\ref{sub:hm-garantien} auf allen \HmEpsAkzeptanzInstanzen{} Instanzen
dieselben Kapazitäten wie das exakte Konfigurationsprogramm und liefert auf
\HmEpsSchlechterRef{} von ihnen dennoch einen größeren Makespan, weil er einen Bin über die
akzeptierte Kapazität hinaus füllen darf.

Die Laufzeiten ordnen sich ebenfalls anders, als es die Güte nahelegt.
Das exakte Konfigurationsprogramm liegt mit \HmZeitMedianggip{}~ms im Median vor dem
$\mathrm{OPT}+1$-Algorithmus mit \HmZeitMedianjansen{}~ms, dessen innere Binärsuche mehrere
Modelle je Kapazität löst.
Wächst die Anzahl der Aufgaben um den Faktor \HmSkalNFaktor, so steigt die Laufzeit von
MULTIFIT+FFD um das \HmSkalFaktorffd-fache, die der High-Multiplicity-Verfahren dagegen nur um
das \HmSkalFaktorMin- bis \HmSkalFaktorMax-fache.

\paragraph{Grenzen der Aussagen.}
Aus den Messungen folgt keine Schranke.
Dass auf \EvalInstanzen{} Instanzen kein Verhältnis über $13/11$ auftritt, ist für
MULTIFIT+FFD mit der bewiesenen Schranke verträglich und begründet für MULTIFIT+MFFD keine.
Ähnliches gilt für das Kollaps-Lemma: Es beschreibt eine Klasse, auf der die Subroutinenwahl
folgenlos bleibt, und gerade der bekannte Extremfall liegt außerhalb, denn die Instanz von
\textcite{friesen1984} hat $\rho = 1{,}65$.
Zur oberen Schranke trägt das Lemma somit nichts bei.

Die zweite Versuchsreihe deckt zudem nur $d \le \HmDMax$ ab, während die fünf
Benchmark-Datensätze im Median zwischen \HmDBenchMedianMin{} und \HmDBenchMedianMax{}
verschiedene Aufgabengrößen enthalten; dort sind die High-Multiplicity-Verfahren nicht gemessen
worden.
Ein Teil der verwendeten Optima stammt aus \textcite{akram2025} und ist nicht nachgerechnet
worden, und sämtliche Laufzeiten gelten für die beschriebene Implementierung und Umgebung.

\paragraph{Offene Fragen.}
Die naheliegendste betrifft MFFD.
Auf der Instanz von \textcite{friesen1984} liefert MULTIFIT+MFFD den kleineren Makespan
(Abschnitt~\ref{sec:friesen}); über das Verhalten im schlechtesten Fall sagt das jedoch nichts,
und ob $r_m(\mathrm{MF}{+}\mathrm{MFFD}) < 13/11$ gilt, bleibt offen.
Dahinter steht die allgemeinere Frage, für welche Subroutinen eine Entsprechung
von~\eqref{eq:ffd-akzeptiert} überhaupt besteht: Für FFD ist sie bewiesen, für MFFD und für das
Konfigurations-LP unbekannt - beim Konfigurations-LP reicht die Rundungsschranke aus
Lemma~\ref{lem:gg-rundung} dafür nicht aus, ein Gegenbeispiel ist damit aber nicht gegeben.
Unbestimmt bleibt ferner das Grenzverhältnis für kleine Maschinenzahlen, denn
Lemma~\ref{lem:untere-schranke} setzt $m \ge 13$ voraus; für $2 \le m < 13$ klafft zwischen
$r_m \ge 1$ aus Bemerkung~\ref{bem:rm-mindestens-eins} und der oberen Schranke $13/11$ eine
Lücke, und nur $r_1 = 1$ ist bestimmt.

Neu entstanden ist die Frage nach der Reichweite der $\varepsilon$-dualen Konstruktion.
Über die von \textcite{jansen_solisoba2010} übernommene Wahl $\varepsilon = 1/(2d-1)$ wird ihre
Garantie mit wachsender Anzahl verschiedener Aufgabengrößen schärfer, die Laufzeit des
Verfahrens nach Abschnitt~\ref{sub:jansen-laufzeit} jedoch doppelt exponentiell in $d$.
Lemma~\ref{lem:jansen-dual} selbst verlangt diese Kopplung nicht; sie begrenzt die Zahl der
großen Konfigurationen und damit die Größe des Programms.
Am schwächsten ist die Garantie demnach dort, wo sie sich berechnen lässt.
Die in Abschnitt~\ref{sub:gg-modell} beschriebene Spaltengenerierung hält das
Konfigurationsmodell auch bei wachsendem $d$ handhabbar; implementiert ist sie hier nicht, und
auf das ganzzahlige Programm des $\varepsilon$-dualen Verfahrens ist sie nicht übertragen
worden.
Zu klären wäre schließlich, ob sich die zertifizierte Ablehnung des Konfigurations-LPs mit einer
schnellen Subroutine verbinden lässt, welche das Suchintervall zwar nicht beweisbar eingrenzt,
dafür aber deutlich schneller entscheidet (Abschnitt~\ref{sub:hm-laufzeit}).
